% LSPV-4 REPORT::READFLOOR1::2026-10-12T18:55:44Z
% !TEX program = xelatex
\documentclass[11pt]{article}
\usepackage{fontspec}
\usepackage[margin=0.86in]{geometry}
\usepackage{amsmath,amssymb,amsthm,mathtools,booktabs,longtable,tabularx,array,enumitem,xurl,hyperref,microtype}
\hypersetup{colorlinks=true,linkcolor=blue,urlcolor=blue,citecolor=blue}
\newtheorem{theorem}{Theorem}
\newtheorem{corollary}{Corollary}
\newtheorem{definition}{Definition}
\newcommand{\Core}{\mathcal C_{2,2}^{\pm}}
\newcommand{\EXACT}{\textbf{PROVED/EXACT}}
\newcommand{\KNOWN}{\textbf{KNOWN THEOREM}}
\newcommand{\COMP}{\textbf{COMPUTATIONAL}}
\newcommand{\CAND}{\textbf{CANDIDATE}}
\newcommand{\OPEN}{\textbf{OPEN}}
\newcommand{\RET}{\textbf{FALSIFIED/RETIRED}}
\newcommand{\same}{\textsc{same\_realization}}

\title{READFLOOR1: A near-linear clean-core probe floor\\
is universal but not a hardness bridge}
\author{LSPV-4 / UNIVERSAL\_CORE\_BRIDGE}
\date{12 October 2026, 18:55 UTC}

\begin{document}
\maketitle

\begin{center}
\fbox{\parbox{0.94\linewidth}{
\textbf{Executive status.} P versus NP remains \OPEN.  On the fixed
duplicate-free signed Monotone \(3\)SAT-\((2,2)\) core, every deterministic
decider has an explicit satisfiable length-\(N_k\) input on which it must read
at least \(k=\Omega(N_k/\log N_k)\) distinct input positions.  This is
\EXACT{} in the stated random-access bit model.  The same block adversary
holds for linear-time Horn/2SAT, affine, and bounded-treewidth languages, and
READCEIL1 gives the matching universal ceiling \(N_k\).  Therefore
\texttt{INPUT-PROBE-FLOOR-AS-DECISION-SELECTION-BRIDGE1} is \RET.  The
post-read computation needed to select the SAT bit remains \OPEN.}}
\end{center}

\section{Claim card, predecessor, and obstacles}
\begin{description}[leftmargin=0cm,labelwidth=0cm]
\item[Claim ID.] \texttt{CLEAN-CORE-PRODUCT-READ-FLOOR1}.
\item[Input family.] \(\Core\), the canonical indexed encodings of
duplicate-free proper signed Monotone \(3\)SAT-\((2,2)\).
\item[Algorithm quantifier.] Every deterministic machine correct on every
valid core encoding; uniformity and a polynomial clock are not needed for the
probe lemma.
\item[Model and cost.] Random-access read-only input bits, arbitrary charged
work tapes, and one output bit.  Every distinct queried input position costs
one; all other time is still charged but is not lower-bounded by this lemma.
\item[Two obstacles.] First, force a resource for every decider rather than a
chosen representation.  Second, make the forced resource strong enough to
contradict polynomial time.  The theorem solves the first only at a near-linear
probe scale and the easy controls falsify the second.
\end{description}

The cyclic predecessor is LSPV-3.  Its newest report,
\path{research/agents/lspv-3/2026-10-13_0236_LSPV-3_REVUNCOMP1.tex}, was read
back from GitHub with the required header and final sentinel.  No recovery was
needed.  The earlier MODELBRIDGE-0022 repair is also already canonical at
\path{research/agents/lspv-4/2026-10-12_1002_LSPV-4_MODELBRIDGE-0022.tex}; it
was not duplicated or reconstructed.

\section{Exact clean blocks and encoding}
Let \(S\) have variables \(x_1,\ldots,x_6\), positive supports
\[
 123,\quad145,\quad246,\quad356,
\]
and the same four supports as negative clauses.  Every variable occurs twice
positively and twice negatively, all clauses are distinct and proper, and
\((x_1,\ldots,x_6)=(0,0,1,1,0,0)\) satisfies \(S\).  Hence \(S\in\Core\).

Döcker constructs an unsatisfiable \(U\in\Core\) with 198 variables and 264
clauses~\cite[Proposition~1]{Doecker19}.  Let \(B_1\) be the disjoint union of
33 copies of \(S\), and let \(B_0=U\).  Both have 198 variables and 264
clauses; \(B_1\) is satisfiable and \(B_0\) is not.  Disjoint union preserves
properness, clause uniqueness, monotonicity, and exactly two occurrences of
each sign.

For \(k\ge1\), allocate \(k\) consecutive 198-variable intervals.  Define
\[
 F_k^*=B_1^{\sqcup k},\qquad
 F_k^{(j)}=B_1^{\sqcup(j-1)}\sqcup B_0\sqcup
 B_1^{\sqcup(k-j)}.                                      \tag{1}
\]
Thus \(F_k^*\) is satisfiable while every \(F_k^{(j)}\) is unsatisfiable.

The declared encoding writes the common header
\((n,m)=(198k,264k)\), followed by the \(k\) block payloads in interval order;
clauses inside a block use one fixed canonical order.  With
\(w_k=\lceil\log_2(198k+1)\rceil\), every signed literal uses one sign bit and
\(w_k\) identifier bits.  Consequently
\[
 N_k=h(k)+792k(1+w_k)+O(k)=\Theta(k\log k),                \tag{2}
\]
where the header \(h(k)=O(\log k)\).  All strings in (1) have exactly the same
length and header.

\section{Universal product probe theorem}
For each \(j\), let \(D_j\) be the nonempty set of payload positions where
\(\operatorname{enc}(F_k^*)\) and
\(\operatorname{enc}(F_k^{(j)})\) differ.  The fixed block layout makes the
sets \(D_1,\ldots,D_k\) pairwise disjoint.

\begin{theorem}[\texttt{CLEAN-CORE-PRODUCT-READ-FLOOR1}]
For every deterministic decider \(A\) correct on \(\Core\) and every
\(k\ge1\), the run of \(A\) on \(\operatorname{enc}(F_k^*)\) queries at least
one position in every \(D_j\).  Hence it makes at least
\[
 k=\Omega(N_k/\log N_k)                                  \tag{3}
\]
distinct input probes and at least that much time.
\end{theorem}

\begin{proof}
Fix \(j\).  If the realized run on \(F_k^*\) never queries \(D_j\), then
\(F_k^*\) and \(F_k^{(j)}\) agree at every queried address.  Induction on the
deterministic transitions gives the same internal states, next addresses,
read values, halt time, and output on both inputs.  This contradicts their
opposite satisfiability labels.  Thus the run queries every \(D_j\).  Since
the \(D_j\)'s are disjoint, it queries at least \(k\) distinct positions.
Equation (2) yields (3).
\end{proof}

The quantifiers are worth making explicit:
\[
 \forall A\in\operatorname{Dec}(\Core)\ \forall k\ge1:\quad
 q_A(F_k^*)\ge k.                                         \tag{4}
\]
There is no algorithm-dependent construction, hidden exponent, or witness
output.  The proof charges one realized run.  Yet (3) is at most near-linear,
while READCEIL1 gives \(q_A(F)\le |\operatorname{enc}(F)|\) for every input.

\section{Compatibility semantics and why the bridge fails}
For any partial assignment \(\alpha\), retain the semantic object
\[
 H(F,\alpha)=\{\beta:\alpha\cup\beta\models F\}.           \tag{5}
\]
Equation (1) ensures \(H(F_k^*,\varnothing)\ne\varnothing\) and
\(H(F_k^{(j)},\varnothing)=\varnothing\).  The theorem forces enough source
observation to distinguish these emptiness cases.  It does not force the
machine to materialize a member of (5), a residual table, a decomposition, or
any compatibility representation.  After all required positions have been
read, the computation choosing the one-bit label may use Gaussian elimination,
dynamic programming, a compact circuit, recomputation, or any other uniform
method.  No superpolynomial lower bound follows from (3).

Assuming \(\Core\in\mathsf P\), a polynomial decider is fully consistent with
the floor: scanning the whole input costs \(O(N_k)\), after which the assumed
algorithm runs in \(N_k^{K_A}\).  Banning that branch would merely assume the
desired separation.  Thus the surviving arrow is
\[
 \Core\in\mathsf P
 \Longrightarrow \text{near-linear source probes}
 \xRightarrow[\text{not proved}]{\OPEN}
 \text{superpolynomial post-read decision work}
 \Longrightarrow\bot .                                   \tag{6}
\]

\section{Hostile controls and representation audit}
The same indistinguishability proof is present in easy languages.  For each
block, use the padded Horn/2CNF pair
\[
 (x\vee x)\wedge(x\vee\neg x)
 \quad\hbox{versus}\quad
 (x\vee x)\wedge(\neg x\vee\neg x),                      \tag{7}
\]
or the affine pair \(x=0,x=0\) versus \(x=0,x=1\).  A product decider must
touch one differing region per block on the all-satisfiable product, although
both families are solvable in polynomial time (indeed linearly for these
products).  The 2CNF pair also has bounded primal treewidth.  Therefore the
probe floor detects input sensitivity, not NP hardness.

COPY, DEAD, DELAY, INTERP, RECOMP, and equivalent wrappers do not evade the
indistinguishability statement for their actual read trace; delays can add
arbitrary work without strengthening the lower bound.  Redundant clauses can
be added symmetrically.  An alternate fixed block code preserves the
\(\Omega(k)\) statement, but padding by ignored bits can make its expression in
total encoded length arbitrarily weaker.  Hence (3) is encoding-aware rather
than a representation-invariant superlinear resource.

\section{Bounded reproducible experiment}
The script
\path{experiments/review/clean_core_product_read_floor.py} writes
\path{results/raw/review/2026-10-13_clean_core_product_read_floor.json}.  One
invocation completed in 0.4 seconds; no chunking was needed.  It made nine
top-level checks with zero failures.  It verified the six-variable clean SAT
block and the arithmetic \(33\cdot6=198\), \(33\cdot8=264\); independently
searched a small equal-size proper 3-CNF SAT/UNSAT pair; exhaustively checked
its truth tables; and checked disjoint difference sets for
\(k=1,2,4,8,16,32,64,128\).  The largest synthetic payload was 38,016 bits.
It also records Horn/2SAT, affine, wrapper, and padding controls.

The script SHA-256 is
\path{80d444be91c9e7e4a1655291c2252cc01c3a243c5ef1648ff89b177e4c8aae47};
the JSON SHA-256 is
\path{c95f87e36a61d5b7705d331df31e29241eb022de8c8af8d3664b5a78c9a4c884}.
The published 198-variable unsatisfiable block is used symbolically from the
cited theorem and is not reconstructed by the harness.  The finite checks do
not establish any asymptotic runtime claim.

\section{Faithful main4/main5 and peer reconciliation}
The actual \path{main4.tex} remains the ARC achievable-region manuscript and
\path{main5.tex} remains the WRAPPER construction manuscript.  Neither says
that every uniform SAT decider realizes its DAG, proof, compatibility object,
or wrapper.  Their all-input acquisition and P-versus-NP arrows remain open.

The newest MAIN REPPROMISE1 proves the representation/promise endpoint
trilemma and leaves \texttt{DECISION-SELECTION-WORK1} open.  HIERFIRE1 retains
arbitrary solver exponents; SUCTYPE1 separates explicit from succinct return
types; REVUNCOMP1 defeats universal erasure charges; READCEIL1 proves the
linear read ceiling.  READFLOOR1 adds the complementary clean-core lower floor
but, because (7) has the same floor, narrows rather than closes the decision-
selection target.  The older BLACKBOX-QUERY-SCALE-DICHOTOMY1 already closes
the subpolynomial versus positive-scale black-box query split; that route was
not renamed or reproved here.

\section{Literature and transfer classification}
Döcker's explicit unsatisfiable core and NP-completeness theorem are
\textbf{DIRECTLY TRANSFERABLE} to the blocks above~\cite{Doecker19}.  The
Rivest--Vuillemin evasiveness method is a primary decision-tree precedent for
adversarial indistinguishability~\cite{RV76}: \textbf{DIRECTLY TRANSFERABLE}
only in the input-query model.  It does not turn a probe bound into an
unrestricted RAM/Turing time lower bound.

Circuit, branching-program, proof-system, VLSI, communication-cut, streaming,
and time--space lower bounds are \textbf{FAILED TRANSFERS} without a compiler
placing the same solver realization in the cited model with all costs charged.
Control-theoretic observability is a \textbf{USEFUL STRUCTURAL ANALOGY} for
distinguishing the block substitutions, but it supplies no computation lower
bound after the input is observed.

The theorem and its easy controls relativize; they do not escape the
Baker--Gill--Solovay barrier~\cite{BGS75}.  They also algebrize in the relevant
sense and provide no Aaronson--Wigderson escape~\cite{AW09}.  No large useful
constructive circuit property is produced, so Natural Proofs is not directly
triggered~\cite{RR97}.  These audits are barriers, not premises in the proof.

\section{Post-quantum technique ledger}
\begin{longtable}{p{0.29\linewidth}p{0.64\linewidth}}
\toprule Technique & Classification and consequence\\ \midrule\endhead
Schur squares, shortening, puncturing, dualization & \textbf{FAILED TRANSFER}:
the queried coordinates do not form a forced linear code representation.\\
Quadratic relations, Hilbert functions, syzygies & \textbf{FAILED TRANSFER}
without a forced graded algebra; hidden structure could instead make the
input easier.\\
Lossiness/collision entropy & \textbf{EXACT OBJECT/THEOREM TRANSFER} for the
map from inputs to read certificates: an unread \(D_j\) would create a
cross-label collision.  Eliminating such collisions costs only one probe per
block.\\
Worst-case/average-case architecture & \textbf{ANALOGY/CANDIDATE}: an external
extractor would need a new distributional theorem and a fully charged
same-realization reduction.\\
Special soundness / conflicting openings & \textbf{FAILED TRANSFER}: the two
formulas are different statements, not two openings of one commitment.\\
Hash/random oracle & \textbf{FAILED TRANSFER} absent unconditional removal of
the assumption.\\ \bottomrule
\end{longtable}

\section{CORPUS REVIEWED}
This run surveyed the GitHub repository tree; read root \path{main.tex},
\path{AGENTS.md}, \path{CODEX.md}, the PNP skill and linked protocol/mode
files, all current \path{tasks/}, actual \path{main4.tex}/\path{main5.tex},
newest verified MAIN REPPROMISE1, and newest LSPV reports HIERFIRE1, SUCTYPE1,
REVUNCOMP1, and READCEIL1.  It inspected materially relevant historical MAIN
snapshots, the black-box query-scale dichotomy, readback/certificate reports,
and the HEC, VSC-HEC, G-PPC-HEC, CRS-PPC, CCSPC, RCCN, and Search-G branches;
governance and cited experiment/result artifacts were also surveyed.  Search
covered product adversaries, input probes, block sensitivity, query scales,
wrappers, padding, core encodings, relativization, algebrization, and current
primary literature.

\section{DUPLICATION / CONTRADICTION AUDIT}
The transcript-indistinguishability argument is elementary decision-tree
methodology, not claimed as a new field theorem.  The new corpus contribution
is the exact equal-size \(\Core\) instantiation, its quantifiers, and the
floor/ceiling/easy-control reconciliation.  DELAYOR1 and
BLACKBOX-POSITIVE-SCALE-GATE1 already handle arbitrary query-length exponents;
BLACKBOX-QUERY-SCALE-DICHOTOMY1 already unifies those regimes.  This report
does not duplicate them.  No main4/main5 assertion is strengthened.

\section{WHAT CHANGED RELATIVE TO ROOT MAIN}
The root previously recorded only the universal read-certificate theorem and
ceiling.  This run adds
\[
 \Omega(N/\log N)\le \max_{F\in\Core,|F|=N}q_A(F)\le N
\]
along the explicit product subsequence \(N=N_k\), for every deterministic
core decider.  It then falsifies the lower floor as a hardness hinge using
easy block products and narrows the open target to post-read decision work or
another obligation strictly stronger than source sensitivity.

\section{ADOPTED / REJECTED / RETIRED PRIOR COMPONENTS}
\textbf{Adopted:} the corrected signed \((2,2)\) core, canonical indexed
encoding, READCEIL1 transition induction, one-bit output, full accounting,
\same{}, arbitrary solver exponents, and the explicit 198-variable published
unsatisfiable core.

\textbf{Rejected as universal transfers:} semantic multiplicity to runtime;
probe count to post-read computation; easy product sensitivity to NP hardness;
and circuit, proof, communication, VLSI, streaming, or service-interface
lower bounds to arbitrary computation.

\textbf{FALSIFIED/RETIRED:}
\texttt{INPUT-PROBE-FLOOR-AS-DECISION-SELECTION-BRIDGE1}; previously retired
read-certificate cardinality, trace materialization, generic erasure, and
black-box scale routes remain retired.

\section{STALE OR OVERSTATED CLAIMS FOUND}
It is stale to say only that a decider \emph{may} read the full source: the
product theorem forces it to touch every one of \(k\) disjoint difference
regions on the explicit all-satisfiable input.  It is nevertheless overstated
to call that source sensitivity a SAT-hard computation lower bound, since the
identical adversary occurs in (7) and affine systems.  A semantic-state count
still gives no runtime bound.

\section{Strict status and handoff}
\begin{center}
\begin{tabularx}{0.98\linewidth}{Xp{0.22\linewidth}X}
\toprule Item & Status & Reason\\ \midrule
Clean-core product probe floor & \EXACT & Deterministic transcript proof\\
Published \(198/264\) unsatisfiable clean block & \KNOWN & Döcker Prop.~1\\
Finite block/encoding/control audit & \COMP & Nine checks, zero failures\\
Probe floor as decision-selection bridge & \RET & Easy P controls share it\\
Post-read decision-selection work & \OPEN & No general-machine lower bound\\
Model-matched unconditional contradiction & \OPEN & Transfer absent\\
P versus NP & \OPEN & Load-bearing arrow open\\ \bottomrule
\end{tabularx}
\end{center}

LSPV-1 should test post-read state obligations against arbitrary solver
exponents.  LSPV-2 should search for a relational observable not already
present in easy block products.  LSPV-3 should apply recoding, padding,
affine/Horn/treewidth controls and equivalent wrappers.  LSPV-5 should record
the clean-core floor/ceiling band but keep \texttt{DECISION-SELECTION-WORK1}
and the model-matched contradiction open.

\textbf{Highest-value next target.} Define an externally testable one-bit
decision obligation that remains nontrivial after all source positions are
available, fails on at least the Horn/2SAT and affine controls, is invariant
under equivalent wrappers and alternate encodings, and carries a
solver-independent superpolynomial lower bound in the same general machine
model.  No such obligation is presently known.

\begin{thebibliography}{99}
\bibitem{Doecker19} J. Döcker, ``Monotone 3-Sat-(2,2) is NP-complete,''
arXiv:1912.08032, 2019, \url{https://arxiv.org/abs/1912.08032}.
\bibitem{RV76} R. L. Rivest and J. Vuillemin, ``On recognizing graph
properties from adjacency matrices,'' \emph{Theoretical Computer Science}
3(3), 371--384, 1976,
\url{https://doi.org/10.1016/0304-3975(76)90053-0}.
\bibitem{BGS75} T. Baker, J. Gill, and R. Solovay, ``Relativizations of the
P=?NP question,'' \emph{SIAM J. Comput.} 4(4), 431--442, 1975,
\url{https://doi.org/10.1137/0204037}.
\bibitem{AW09} S. Aaronson and A. Wigderson, ``Algebrization: a new barrier in
complexity theory,'' \emph{ACM TOCT} 1(1), 2009,
\url{https://doi.org/10.1145/1490270.1490272}.
\bibitem{RR97} A. Razborov and S. Rudich, ``Natural proofs,'' \emph{JCSS}
55(1), 24--35, 1997,
\url{https://doi.org/10.1006/jcss.1997.1494}.
\end{thebibliography}

% END_OF_VERIFIED_REPORT::LSPV-4::READFLOOR1
\end{document}
